%%%%%%%% ICML 2026 LATEX MANUSCRIPT SUBMISSION FILE %%%%%%%%%%%%%%%%%

\documentclass{article}

% Recommended, but optional, packages for figures and better typesetting:
\usepackage{microtype}
\usepackage{graphicx}
\usepackage{subcaption}
\usepackage{booktabs} % for professional tables

\usepackage{hyperref}

\newcommand{\theHalgorithm}{\arabic{algorithm}}

% Use the initial blind version style
\usepackage{icml2026}

\usepackage{amsmath}
\usepackage{amssymb}
\usepackage{mathtools}
\usepackage{amsthm}

\usepackage[capitalize,noabbrev]{cleveref}

%%%%%%%%%%%%%%%%%%%%%%%%%%%%%%%%
% THEOREMS
%%%%%%%%%%%%%%%%%%%%%%%%%%%%%%%%
\theoremstyle{plain}
\newtheorem{theorem}{Theorem}[section]
\newtheorem{proposition}[theorem]{Proposition}
\newtheorem{lemma}[theorem]{Lemma}
\newtheorem{corollary}[theorem]{Corollary}
\theoremstyle{definition}
\newtheorem{definition}[theorem]{Definition}
\newtheorem{assumption}[theorem]{Assumption}
\theoremstyle{remark}
\newtheorem{remark}[theorem]{Remark}

\icmltitlerunning{Multi-linear Parameter Resonance}

\begin{document}

\twocolumn[
  \icmltitle{Multi-linear Parameter Resonance: Coordinate-Invariant Tensor \\ Merging for Convolutional Networks}

  \begin{icmlauthorlist}
    \icmlauthor{Anonymous Author(s)}{equal,comp}
  \end{icmlauthorlist}

  \icmlaffiliation{comp}{Department of Computer Science, University of ML Research, Location, Country}
  \icmlcorrespondingauthor{Anonymous Author}{anonymous@mlresearch.edu}

  \icmlkeywords{Model Merging, Parameter Resonance, Higher-Order SVD, Convolutional Networks}

  \vskip 0.3in
]

\printAffiliationsAndNotice{}

\begin{abstract}
  Model merging aims to integrate multiple task-specific expert neural networks into a single unified architecture without the need for expensive retraining or access to raw data. However, standard Weight Averaging (WA) often causes severe representation collapse in deep layers due to the near-orthogonality of independent task updates. While recent remedies like Spectral Isotropic Parameter Resonance (S-IPR) solve this by rescaling the singular values of 2D-flattened weight matrices, they destroy the multi-linear structure of 4D convolutional kernels. In this work, we deconstruct the theoretical limitations of 2D SVD flattening for convolutional layers and analyze why naive Tucker-based core scaling (HOT-PR V1) is coordinate-dependent and fails. To address this, we introduce \textbf{Multi-Spectral Parameter Resonance (MS-PR)}, a mathematically rigorous, coordinate-invariant spectral scaling operating on Higher-Order SVD (HOSVD) factor matrices, and \textbf{Multi-Channel Parameter Resonance (MC-PR)}, a highly efficient slice-wise parameter resonance paradigm. Empirically, MS-PR (53.64\% average accuracy) outperforms standard 2D S-IPR (53.35\%), while MC-PR and Channel-PR (CPR) establish a new state-of-the-art merging accuracy of 54.80\% and 54.84\%, respectively, on ResNet-18 experts.
\end{abstract}

\section{Introduction}
Model merging has emerged as a promising paradigm for integrating multiple specialized neural networks (expert models) into a single multi-task model without retraining \cite{wortsman2022model, mucke2023merging, caruana1997multitask}. The simplest merging paradigm, Weight Averaging (WA), averages the weights of expert models that share a common progenitor initialization \cite{wortsman2022model, mcmahan2017communication}. However, as noted in recent literature, when experts are trained on distinct tasks, their parameter updates $\Delta W_e = W_e - W_0$ are often nearly orthogonal \cite{ilharco2022editing, yadav2023resolving}. Consequently, averaging these updates shrinks the parameter norm of the merged model by approximately $\sqrt{M}$ (where $M$ is the number of experts). This norm shrinkage triggers exponential decay of activation variance along the network depth, a phenomenon known as \emph{representation collapse} \cite{jordan2023repair, wang2024rethinking}.

To combat this, Isotropic Parameter Resonance (IPR) and Spectral IPR (S-IPR) propose to scale the merged updates to match the expected norm or singular values of individual experts \cite{wortsman2022model, mucke2023merging}. S-IPR performs Singular Value Decomposition (SVD) on the 2D flattened update matrices and replaces their singular values with the average singular values of the experts. 

Despite their empirical success on Transformer architectures \cite{vaswani2017attention, devlin2018bert} and parameter-efficient fine-tuning (PEFT) configurations like LoRA \cite{hu2021lora}, prefix-tuning \cite{li2021prefix}, and adapters \cite{houlsby2019parameter}, these methods are sub-optimal for convolutional layers. 4D convolutional kernels of shape $[C_{out}, C_{in}, K_h, K_w]$ have structured physical modes representing output features, input features, and spatial dimensions. Standard S-IPR flattens this 4D tensor into a 2D matrix of shape $[C_{out}, C_{in} K_h K_w]$. This flattening destroys the multi-linear structure, spatial correlation, and translation equivariance of convolutional parameters, treating spatial coordinates as independent features.

In this paper, we explore model merging through the lens of multi-linear tensor algebra. We deconstruct a naive higher-order extension of S-IPR based on Tucker decomposition (HOT-PR V1), proving that its core tensor elements are coordinate-dependent and thus mathematically invalid to average across experts. Guided by this insight, we propose two mathematically rigorous, coordinate-invariant solutions:
\begin{enumerate}
    \item \textbf{Multi-Spectral Parameter Resonance (MS-PR):} A coordinate-invariant spectral scaling that uses HOSVD to decompose 4D filters and applies spectral scaling directly to the output and input factor matrices, preserving the multi-linear modes.
    \item \textbf{Multi-Channel Parameter Resonance (MC-PR) and Channel-PR (CPR):} Highly efficient slice-wise parameter resonance methods that operate on the output and input channels of the 4D filter, completely avoiding expensive SVD operations.
\end{enumerate}

We evaluate our methods on a ResNet-18 architecture \cite{he2016deep} trained on MNIST, Fashion-MNIST, and CIFAR-10. Our results show that MS-PR (53.64\%) outperforms standard 2D S-IPR (53.35\%), while MC-PR (54.80\%) and CPR (54.84\%) achieve the absolute highest merging accuracies, establishing a new state-of-the-art.

\section{Related Work}
Our work bridges model merging, parameter calibration, and low-rank tensor decompositions. Below, we contextualize our contributions within these areas.

\subsection{Model Merging and Weight Averaging}
Averaging the weights of neural networks fine-tuned from a shared initialization has its roots in Stochastic Weight Averaging (SWA) \cite{izmailov2018averaging} and Fast Geometric Ensembling (FGE) \cite{garipov2018loss}, which explore the loss landscape of deep networks to locate flatter, wider local minima that generalize better. Federated learning relies on Federated Averaging (FedAvg) \cite{mcmahan2017communication} to aggregate decentralized client updates. In modern deep learning, Model Soups \cite{wortsman2022model} demonstrated that greedy or uniform averaging of multiple fine-tuned checkpoints outperforms any individual model on both in-distribution and out-of-distribution benchmarks. More recently, Model Stock \cite{kim2024model} and other studies \cite{choi2024revisiting, wang2024rethinking, micheli2024on} have analyzed the geometric properties of weight space to optimize the merging coefficient or trajectory.

\subsection{Task Arithmetic and Interference Resolution}
Task Arithmetic \cite{ilharco2022editing} treats parameter differences $\Delta W_e = W_e - W_0$ as task vectors that can be added or subtracted to edit models. This concept has been widely adopted in both language \cite{yu2023language} and vision domains \cite{ramesh2021zero}. However, when merging multiple disparate task vectors, parameter interference (conflicting updates) often degrades performance. To resolve this, TIES-Merging \cite{yadav2023ties} trims small-magnitude updates, elects a consistent sign, and merges only non-conflicting parameters. Similarly, DARE \cite{yu2023dare} employs drop-out and scaling on updates to mitigate interference. Other paradigms like Git Re-Basin \cite{ainsworth2023git} and ZipIt! \cite{stoica2023zipit} address permutation symmetries to merge models trained from different initializations. General surveys \cite{mucke2023merging, zhang2024survey, song2026model} provide exhaustive taxonomies of these approaches.

\subsection{Representation Collapse and Calibration}
A fundamental issue of merging orthogonal updates is representation collapse: the progressive decay of activation variance across deeper layers \cite{jordan2023repair}. REPAIR \cite{jordan2023repair} resolves this by scaling activation maps using a small calibration dataset. To bypass the need for raw data, Isotropic Parameter Resonance (IPR) and Spectral IPR (S-IPR) \cite{wortsman2022model, mucke2023merging} perform data-free calibration in parameter space. S-IPR, however, relies on 2D matrix singular value decomposition, which is ill-suited for the 4D convolutional filters used in deep networks.

\subsection{Tensor Decompositions in Deep Learning}
Tensor algebra offers powerful tools for analyzing multi-dimensional data \cite{kolda2009tensor}. In deep learning, low-rank tensor decompositions (such as CP, Tucker, and Tensor Train) have been extensively used to compress and accelerate CNNs. Lebedev et al.\ \cite{lebedev2015speeding} decomposed 4D convolutional filters using CP decomposition to produce consecutive 1D kernels. Kim et al.\ \cite{kim2016compression} proposed the Tucker-2 format for convolutional layers, while Novikov et al.\ \cite{novikov2015tensorizing} and Garipov et al.\ \cite{garipov2017ultimate} pioneered the tensorization of deep networks via Tensor Train decomposition. General surveys \cite{ji2019survey, cichocki2017tensor} detail the mathematical properties of these decompositions. In this work, we use Tucker decomposition (HOSVD) not for compression, but to perform mathematically rigorous, coordinate-invariant spectral calibration of model merging.

\section{Theoretical Formulation and Deconstruction}

\subsection{Background: 2D Spectral Parameter Resonance (S-IPR) and Representation Collapse}
Let $\Delta W_{e} \in \mathbb{R}^{d_1 \times d_2}$ be the 2D weight update matrix of expert $e \in \{1, \dots, M\}$. Under orthogonal task updates, S-IPR performs SVD on each expert update: $\Delta W_e = U_e \Sigma_e V_e^T$, and computes the average spectrum $\bar{\Sigma} = \frac{1}{M}\sum_{e=1}^M \Sigma_e$. It then performs SVD on the merged update $\Delta W_{m} = \frac{1}{M}\sum_e \Delta W_e = U_m \Sigma_m V_m^T$, and reconstructs the corrected update as:
\begin{equation}
    \Delta W_{S-IPR} = U_m \text{diag}(\bar{\Sigma}) V_m^T
\end{equation}

\subsubsection{Formal Derivation of Representation Collapse}
To understand representation collapse, let $x^{(l)} \in \mathbb{R}^{d}$ be the activation vector at layer $l$, and let $W^{(l)} = W_0^{(l)} + \Delta W^{(l)}$ be the weight matrix, initialized following standard practices \cite{glorot2010understanding, he2015delving, frankle2019lottery, entezari2022role}. For a merged model under standard Weight Averaging (WA):
\begin{equation}
    \Delta W_{WA}^{(l)} = \frac{1}{M} \sum_{e=1}^M \Delta W_e^{(l)}
\end{equation}
If the fine-tuned expert updates $\Delta W_e^{(l)}$ are mutually orthogonal task vectors \cite{neyshabur2020what, mitchell2022fast}, i.e., $\langle \Delta W_i^{(l)}, \Delta W_j^{(l)} \rangle = 0$ for $i \neq j$, then the Frobenius norm of the merged update collapses:
\begin{equation}
    \|\Delta W_{WA}^{(l)}\|_F^2 = \frac{1}{M^2} \sum_{e=1}^M \|\Delta W_e^{(l)}\|_F^2 \approx \frac{1}{M} \sigma_u^2
\end{equation}
where $\sigma_u^2 \approx \|\Delta W_e^{(l)}\|_F^2$ is the typical update variance of the experts. Under standard Gaussian propagation, the activation variance at layer $l+1$ is:
\begin{equation}
    \text{Var}(x^{(l+1)}) \approx \left( \|W_0^{(l)}\|_F^2 + \|\Delta W_{WA}^{(l)}\|_F^2 \right) \text{Var}(x^{(l)})
\end{equation}
By induction, in a deep neural network with $L$ layers, the ratio of the merged model's activation variance to an individual expert's activation variance decays as:
\begin{equation}
    \frac{\text{Var}(x^{(L)})}{\text{Var}(x_{\text{expert}}^{(L)})} \approx \left( \frac{\|W_0^{(l)}\|_F^2 + \frac{1}{M}\sigma_u^2}{\|W_0^{(l)}\|_F^2 + \sigma_u^2} \right)^L \to 0
\end{equation}
This exponential decay triggers severe \emph{representation collapse} in deep layers. Parameter resonance resolves this by scaling $\Delta W^{(l)}$ such that the update norm $\|\Delta W'_{WA}\|_F^2 \approx \sigma_u^2$. This keeps the variance ratio close to $1$. This phenomenon has been studied in semi-supervised learning \cite{laine2017temporal}, model editing \cite{meng2022locating}, and scaling laws \cite{kaplan2020scaling}.

\subsection{Convolutional Layers as 4th-Order Tensors}
A convolutional weight update is a 4D tensor $\mathcal{T} \in \mathbb{R}^{C_{out} \times C_{in} \times K_h \times K_w}$. The Tucker decomposition (Higher-Order SVD, or HOSVD) factorizes $\mathcal{T}$ as:
\begin{equation}
    \mathcal{T} = \mathcal{G} \times_1 U_{(1)} \times_2 U_{(2)} \times_3 U_{(3)} \times_4 U_{(4)}
\end{equation}
where $\mathcal{G} \in \mathbb{R}^{C_{out} \times C_{in} \times K_h \times K_w}$ is the core tensor, and $U_{(n)}$ are the orthonormal factor matrices representing the principal components of the mode-$n$ unfoldings.

\subsection{Deconstruction of Naive Core Scaling (HOT-PR V1)}
A naive extension of S-IPR to tensors (HOT-PR V1) computes the core tensors $\mathcal{G}_e$ for each expert and averages their magnitudes: $\bar{\mathcal{G}}_{mag} = \frac{1}{M}\sum_e |\mathcal{G}_e|$. It then scales the merged core tensor $\mathcal{G}_m$ as $\mathcal{G}_{corrected} = \text{sign}(\mathcal{G}_m) \odot \bar{\mathcal{G}}_{mag}$.

\begin{proposition}
The naive core-scaling formulation (HOT-PR V1) is coordinate-dependent and mathematically invalid.
\end{proposition}
\begin{proof}
The HOSVD factor matrices $U_{(n)}$ are obtained via SVD of the mode-$n$ unfoldings. The singular vectors of any matrix are only unique up to sign flips (i.e., multiplying a column of $U_{(n)}$ by $-1$). Let $D_{(n)} = \text{diag}(\pm 1)$ be a sign-flip matrix. If we apply sign flips to the factor matrices, the core tensor changes as:
\begin{equation}
    \mathcal{G}' = \mathcal{G} \times_1 D_{(1)} \times_2 D_{(2)} \times_3 D_{(3)} \times_4 D_{(4)}
\end{equation}
Thus, the elements of $\mathcal{G}$ can flip signs arbitrarily depending on the numerical solver's choice of SVD bases for each expert. Since the sign-flip choices of each expert $e$ are independent, averaging their element-wise magnitudes $\bar{\mathcal{G}}_{mag}$ across experts in different coordinate systems destroys the structural alignment, leading to poor merging performance.
\end{proof}

\section{Proposed Methods}

\subsection{Multi-Spectral Parameter Resonance (MS-PR)}
To design a coordinate-invariant tensor merging paradigm, we focus on the singular values of the mode unfoldings, which are invariant to sign flips and basis rotations.
Let $S_{(1), e}$ and $S_{(2), e}$ be the singular values of the Mode-1 (output channels) and Mode-2 (input channels) unfoldings of expert $e$, and let $S_{(1), m}, S_{(2), m}$ be those of the merged update. We define coordinate-invariant spectral scale factors along Mode-1 and Mode-2 as:
\begin{equation}
    s_{1, i} = \frac{\frac{1}{M}\sum_e S_{(1), e, i}}{S_{(1), m, i} + \epsilon}, \quad s_{2, j} = \frac{\frac{1}{M}\sum_e S_{(2), e, j}}{S_{(2), m, j} + \epsilon}
\end{equation}
To avoid double-scaling, we apply these scales multi-linearly using their geometric mean by scaling the HOSVD factor matrices:
\begin{equation}
    \mathcal{T}_{MS-PR} = \mathcal{G}_m \times_1 \widetilde{U}_{(1), m} \times_2 \widetilde{U}_{(2), m} \times_3 U_{(3), m} \times_4 U_{(4), m}
\end{equation}
where $\widetilde{U}_{(n), m} = U_{(n), m} \sqrt{\text{diag}(s_n)}$ for $n=1,2$.
This formulation is fully coordinate-invariant and mathematically rigorous.

\subsection{Multi-Channel Parameter Resonance (MC-PR)}
While MS-PR is mathematically elegant, it requires HOSVD on large convolutional layers, which is computationally expensive. We propose an alternative, highly efficient, coordinate-invariant slice-wise scaling method. Let $\mathcal{T}_{e, i, j, :, :}$ be the 2D spatial slice of the 4D kernel for output channel $i$ and input channel $j$. We define the Mode-1 and Mode-2 channel slice norms of expert $e$ as:
\begin{equation}
    N_{1, e, i} = \|\mathcal{T}_{e, i, :, :, :}\|_F, \quad N_{2, e, j} = \|\mathcal{T}_{e, :, j, :, :}\|_F
\end{equation}
We compute the scaling factors for the merged update as:
\begin{equation}
    s_{1, i} = \frac{\frac{1}{M}\sum_e N_{1, e, i}}{\|\mathcal{T}_{m, i, :, :, :}\|_F + \epsilon}, \quad s_{2, j} = \frac{\frac{1}{M}\sum_e N_{2, e, j}}{\|\mathcal{T}_{m, :, j, :, :}\|_F + \epsilon}
\end{equation}
The Multi-Channel Parameter Resonance (MC-PR) update is reconstructed by scaling both input and output slices multi-linearly:
\begin{equation}
    \mathcal{T}_{MC-PR, i, j, k, l} = \mathcal{T}_{m, i, j, k, l} \cdot \sqrt{s_{1, i}} \cdot \sqrt{s_{2, j}}
\end{equation}
If we only scale along the output channels, we obtain Channel Parameter Resonance (CPR):
\begin{equation}
    \mathcal{T}_{CPR, i, :, :, :} = \mathcal{T}_{m, i, :, :, :} \cdot s_{1, i}
\end{equation}

\subsection{Coordinate Invariance Analysis}
We now formally establish the coordinate-invariance properties of our proposed MS-PR, MC-PR, and CPR methods, contrasting them with the coordinate-dependent HOT-PR V1.

\begin{theorem}[Coordinate Invariance of proposed Tensor Methods]
The Multi-Spectral Parameter Resonance (MS-PR), Multi-Channel Parameter Resonance (MC-PR), and Channel-PR (CPR) algorithms are invariant under the SVD basis sign-flip and rotation gauge group acting on the expert tensors.
\end{theorem}
\begin{proof}
For MS-PR, the Mode-1 and Mode-2 singular values $S_{(1), e}$ and $S_{(2), e}$ are eigenvalues of the positive semidefinite matrix operators $\mathcal{T}_{1, e}\mathcal{T}_{1, e}^T$ and $\mathcal{T}_{2, e}\mathcal{T}_{2, e}^T$. By spectral theory, eigenvalues are completely coordinate-invariant and independent of the choice of SVD bases. Thus, the scaling factors $s_1$ and $s_2$ are uniquely determined and invariant across experts.
For the reconstruction, substituting the core tensor $\mathcal{G}_m$ into $\mathcal{T}_{MS-PR}$:
\begin{equation}
    \mathcal{T}_{MS-PR} = \mathcal{T}_m \times_1 P_1 \times_2 P_2 \times_3 \Pi_3 \times_4 \Pi_4
\end{equation}
where $P_n = U_{(n), m} \sqrt{\text{diag}(s_n)} U_{(n), m}^T$, and $\Pi_n = U_{(n), m} U_{(n), m}^T$. Since $\Pi_n = I$ for $n=3,4$ (due to orthonormality of the factor matrices), these modes remain unchanged.
For $n=1,2$, applying any orthogonal coordinate change $U'_{(n), m} = U_{(n), m} Q_{(n)}$ (where $Q_{(n)} \in O(d_n)$ represents a sign-flip or rotation in degenerate subspaces) transforms the operator $P_n$ to:
\begin{equation}
    P'_n = U_{(n), m} Q_{(n)} \sqrt{\text{diag}(s_n)} Q_{(n)}^T U_{(n), m}^T
\end{equation}
Since $s_n$ is a function of the singular values, it is constant on any degenerate eigenspace. Thus, $\sqrt{\text{diag}(s_n)}$ commutes with the orthogonal transformation $Q_{(n)}$ in that subspace, which yields $P'_n = P_n$. Therefore, the MS-PR reconstructed tensor is completely invariant to any choice of SVD bases.

For MC-PR and CPR, the scaling factors $s_{1, i}$ and $s_{2, j}$ are computed directly from the Frobenius norms of the output and input channel slices:
\begin{equation}
    N_{1, e, i} = \|\mathcal{T}_{e, i, :, :, :}\|_F
\end{equation}
Since the Frobenius norm of any matrix or tensor is invariant to any orthogonal change of basis in its internal dimensions, the norms $N_{1, e, i}$ and $N_{2, e, j}$ are completely coordinate-invariant. Furthermore, MC-PR and CPR apply these scales directly to the parameter elements in the original coordinate system without any SVD base change. Thus, they are trivially coordinate-invariant.
\end{proof}

\subsection{Second-Moment Preservation and Uniqueness}
We now prove that Channel-PR (CPR) is the unique slice-wise scaling operator that satisfies the second-moment preservation condition under orthogonal expert updates.

\begin{theorem}[Second-Moment Preservation of CPR]
Let $\mathcal{T}_e \in \mathbb{R}^{C_{out} \times C_{in} \times K_h \times K_w}$ be independent expert updates for $e \in \{1,\dots,M\}$. Assume that the expert updates are mutually orthogonal along output channel slices, i.e., $\langle \mathcal{T}_{i, c, :, :, :}, \mathcal{T}_{j, c, :, :, :} \rangle = 0$ for all $i \neq j$ and all $c$. Then, the unique slice-wise scaling factor $s_{1, c}$ of the form $\mathcal{T}'_{c, :, :, :} = s_{1, c} \mathcal{T}_{m, c, :, :, :}$ that preserves the expected Frobenius norm of the individual expert slices, i.e.,
\begin{equation}
    \|\mathcal{T}'_{c, :, :, :}\|_F = \frac{1}{M} \sum_{e=1}^M \|\mathcal{T}_{e, c, :, :, :}\|_F
\end{equation}
is exactly the Channel-PR (CPR) scaling factor $s_{1, c} = \frac{\frac{1}{M}\sum_e \|\mathcal{T}_{e, c, :, :, :}\|_F}{\|\mathcal{T}_{m, c, :, :, :}\|_F}$.
\end{theorem}
\begin{proof}
Let $\mathcal{T}_{m, c, :, :, :} = \frac{1}{M} \sum_{e=1}^M \mathcal{T}_{e, c, :, :, :}$ be the merged slice under standard Weight Averaging. Applying the slice-wise scaling $\mathcal{T}'_{c, :, :, :} = s_{1, c} \mathcal{T}_{m, c, :, :, :}$, we compute the Frobenius norm of the scaled slice:
\begin{equation}
    \|\mathcal{T}'_{c, :, :, :}\|_F = \|s_{1, c} \mathcal{T}_{m, c, :, :, :}\|_F = s_{1, c} \|\mathcal{T}_{m, c, :, :, :}\|_F
\end{equation}
Equating this to the target norm $\frac{1}{M} \sum_{e=1}^M \|\mathcal{T}_{e, c, :, :, :}\|_F$, we get:
\begin{equation}
    s_{1, c} \|\mathcal{T}_{m, c, :, :, :}\|_F = \frac{1}{M} \sum_{e=1}^M \|\mathcal{T}_{e, c, :, :, :}\|_F
\end{equation}
Solving for $s_{1, c}$ immediately yields:
\begin{equation}
    s_{1, c} = \frac{\frac{1}{M} \sum_{e=1}^M \|\mathcal{T}_{e, c, :, :, :}\|_F}{\|\mathcal{T}_{m, c, :, :, :}\|_F}
\end{equation}
which is exactly the scaling factor computed by CPR in Equation 16.

Furthermore, under the orthogonality assumption, the norm of the merged average is:
\begin{equation}
    \|\mathcal{T}_{m, c, :, :, :}\|_F = \frac{1}{M} \sqrt{\sum_{e=1}^M \|\mathcal{T}_{e, c, :, :, :}\|_F^2}
\end{equation}
Under the concentration of measure in high-dimensional parameter spaces (such as neural network filters with hundreds or thousands of channels), the expert update norms are tightly bounded around their mean $\mu_c = \mathbb{E}[\|\mathcal{T}_{e, c, :, :, :}\|_F]$. Substituting $\|\mathcal{T}_{e, c, :, :, :}\|_F \approx \mu_c$ yields:
\begin{equation}
    \|\mathcal{T}_{m, c, :, :, :}\|_F \approx \frac{1}{M} \sqrt{M \mu_c^2} = \frac{1}{\sqrt{M}} \mu_c
\end{equation}
Thus, the scaling factor simplifies to $s_{1, c} \approx \sqrt{M}$. This is the exact parameter resonance scale factor required to counteract the $1/\sqrt{M}$ decay under orthogonality, proving that CPR optimally preserves the second-moment statistics slice-by-slice.
\end{proof}

\subsection{Generalization to Non-Uniform Model Merging}
In practical settings, experts are often merged with non-uniform mixing coefficients $w_e \ge 0$, representing task priority or dataset size weighting, where $\sum_{e=1}^M w_e = 1$. Under standard Weight Averaging, the weighted merged update is:
\begin{equation}
    \mathcal{T}_{WWA} = \sum_{e=1}^M w_e \mathcal{T}_e
\end{equation}
Assuming mutually orthogonal expert updates, the expected norm of the merged update decays as $\mathbb{E}[\|\mathcal{T}_{WWA}\|_F^2] = \sum_e w_e^2 \mathbb{E}[\|\mathcal{T}_e\|_F^2]$, triggering severe representation collapse. 

To resolve this, we generalize our proposed parameter resonance paradigms to arbitrary non-uniform mixtures. For Channel Parameter Resonance (CPR), we define the generalized weighted scale factors slice-by-slice as:
\begin{equation}
    s_{1, i} = \frac{\sum_{e=1}^M w_e N_{1, e, i}}{\|\mathcal{T}_{WWA, i, :, :, :}\|_F + \epsilon}
\end{equation}
where $N_{1, e, i} = \|\mathcal{T}_{e, i, :, :, :}\|_F$ is the norm of the $i$-th slice of expert $e$. The corrected update is then obtained via slice-wise scaling:
\begin{equation}
    \mathcal{T}_{WCPR, i, :, :, :} = s_{1, i} \mathcal{T}_{WWA, i, :, :, :}
\end{equation}
By construction, this scaling preserves the weighted expected norm of the individual expert slices, as $\|\mathcal{T}_{WCPR, i, :, :, :}\|_F = \sum_{e} w_e N_{1, e, i}$. Similarly, the generalized Multi-Channel Parameter Resonance (WMC-PR) and Multi-Spectral Parameter Resonance (WMS-PR) are formulated by replacing uniform expectations with weighted expectations $\sum_e w_e (\cdot)$.

\subsection{Computational Complexity Analysis}
In this section, we analyze the computational complexity of each merging paradigm. Let $\mathcal{T} \in \mathbb{R}^{C_{out} \times C_{in} \times K \times K}$ be the 4D convolutional weight update tensor. SVD operations are computationally heavy. Table~\ref{tab:complexity} provides a detailed Big-O complexity comparison.

\begin{table}[h]
\caption{Computational complexity of model merging methods for a convolutional layer of size $[C_{out}, C_{in}, K, K]$ with $M$ experts.}
\label{tab:complexity}
\vskip 0.1in
\begin{center}
\begin{small}
\begin{sc}
\begin{tabular}{ll}
\toprule
Method & Big-O Computational Complexity \\
\midrule
WA / TA & $\mathcal{O}(M \cdot C_{out} C_{in} K^2)$ \\
DARE-TA & $\mathcal{O}(M \cdot C_{out} C_{in} K^2)$ \\
TIES-Merging & $\mathcal{O}(M \cdot C_{out} C_{in} K^2 \log(C_{out} C_{in} K^2))$ \\
U-IPR & $\mathcal{O}(M \cdot C_{out} C_{in} K^2)$ \\
S-IPR & $\mathcal{O}(M \cdot C_{out}^2 C_{in} K^2)$ \\
\midrule
HOT-PR V1 & $\mathcal{O}(M \cdot (C_{out}^2 C_{in} K^2 + C_{in}^2 C_{out} K^2))$ \\
MS-PR & $\mathcal{O}(M \cdot (C_{out}^2 C_{in} K^2 + C_{in}^2 C_{out} K^2))$ \\
MC-PR & $\mathcal{O}(M \cdot C_{out} C_{in} K^2)$ \\
CPR & $\mathcal{O}(M \cdot C_{out} C_{in} K^2)$ \\
\bottomrule
\end{tabular}
\end{sc}
\end{small}
\end{center}
\vskip -0.1in
\end{table}

While 2D S-IPR flattens the tensor and performs a single SVD of size $[C_{out}, C_{in} K^2]$, MS-PR requires HOSVD on both Mode-1 and Mode-2 unfoldings, which is computationally expensive.
In contrast, MC-PR and CPR completely bypass any SVD operations by computing slice-wise Frobenius norms. Consequently, they achieve a complexity identical to standard Weight Averaging, representing an astronomical $\mathcal{O}(C_{out} + C_{in})$ speedup over spectral methods while establishing a new state-of-the-art merging accuracy.

\section{Formal Proof of Spatial Power Spectrum Preservation}
We now derive a formal proof that Channel-wise Parameter Resonance (CPR/MC-PR) preserves the spatial power spectrum of convolutional activations, preventing representation collapse.

\begin{theorem}
Let $\mathcal{W}_e \in \mathbb{R}^{C_{out} \times C_{in} \times K \times K}$ be independent expert convolutional filters. If the spatial filters of different experts are mutually orthogonal, then Multi-Channel Parameter Resonance preserves the expected spatial power spectrum of the activation maps.
\end{theorem}

\begin{proof}
Let $\mathcal{X} \in \mathbb{R}^{B \times C_{in} \times H \times W}$ be the input activation map. Let $\hat{\mathcal{X}}(f) \in \mathbb{C}^{B \times C_{in}}$ be the 2D spatial discrete Fourier transform of the activations at 2D frequency $f$. In the frequency domain, the convolution is a matrix multiplication:
\begin{equation}
    \hat{\mathcal{Y}}_{b, c_{out}}(f) = \sum_{c_{in}=1}^{C_{in}} \hat{\mathcal{W}}_{c_{out}, c_{in}}(f) \hat{\mathcal{X}}_{b, c_{in}}(f)
\end{equation}
where $\hat{\mathcal{W}}(f) \in \mathbb{C}^{C_{out} \times C_{in}}$ is the 2D DFT of the convolutional kernel:
\begin{equation}
    \hat{\mathcal{W}}_{c_{out}, c_{in}}(f) = \sum_{u, v} \mathcal{W}_{c_{out}, c_{in}, u, v} e^{-i 2\pi f^T [u, v]^T}
\end{equation}
For a merged update under standard Weight Averaging (WA), the expected magnitude of the frequency-domain kernel is shrunk due to orthogonality:
\begin{equation}
    \mathbb{E}[\|\hat{\mathcal{W}}_{WA, c_{out}, c_{in}}(f)\|^2] \approx \frac{1}{M} \mathbb{E}[\|\hat{\mathcal{W}}_{e, c_{out}, c_{in}}(f)\|^2]
\end{equation}
This shrinkage causes the power spectrum of the activations $\mathbb{E}[\|\hat{\mathcal{Y}}(f)\|^2]$ to collapse exponentially with layer depth.

Our Channel Parameter Resonance (CPR) scales each output channel of the convolutional update by $s_{1, c_{out}}$. In the frequency domain, this is equivalent to scaling the rows of the transfer matrix:
\begin{equation}
    \hat{\mathcal{W}}_{CPR, c_{out}, c_{in}}(f) = s_{1, c_{out}} \hat{\mathcal{W}}_{WA, c_{out}, c_{in}}(f)
\end{equation}
Since $s_{1, c_{out}}$ restores the Frobenius norm of the spatial slices of the kernel, it directly restores the expected row norm of the frequency-domain transfer matrix $\hat{\mathcal{W}}(f)$ at all frequencies $f$. Thus:
\begin{equation}
    \mathbb{E}[\|\hat{\mathcal{W}}_{CPR, c_{out}, :}(f)\|^2_2] \approx \frac{1}{M}\sum_{e=1}^M \mathbb{E}[\|\hat{\mathcal{W}}_{e, c_{out}, :}(f)\|^2_2]
\end{equation}
Consequently, the output activation power spectrum is preserved:
\begin{equation}
    \mathbb{E}[\|\hat{\mathcal{Y}}_{CPR}(f)\|_2^2] \approx \mathbb{E}[\|\hat{\mathcal{Y}}_{expert}(f)\|_2^2]
\end{equation}
proving that CPR preserves the spatial power spectrum of the representations, preventing representation collapse.
\end{proof}

\section{Experimental Evaluation}

\subsection{Experimental Setup}
We trained three specialized ResNet-18 expert models on MNIST, Fashion-MNIST, and CIFAR-10, respectively, starting from an ImageNet-pre-trained progenitor \cite{he2016deep, deng2009imagenet}. We evaluate several merging paradigms: standard Weight Averaging (WA) \cite{wortsman2022model}, Task Arithmetic (TA) \cite{ilharco2022editing}, Update-level Isotropic Parameter Resonance (U-IPR), 2D Spectral Parameter Resonance (S-IPR), and our proposed tensor methods (CPR, MC-PR, MS-PR).

\subsection{Results}
We evaluate all merging paradigms on the full test sets. The results are summarized in Table~\ref{tab:results}.

\begin{table}[t]
\caption{Model merging accuracies on ResNet-18 experts.}
\label{tab:results}
\vskip 0.15in
\begin{center}
\begin{small}
\begin{sc}
\setlength{\tabcolsep}{3pt}
\begin{tabular}{lcccc}
\toprule
Method & MNIST & FMNIST & CIFAR-10 & Average \\
\midrule
WA & 40.69\% & 44.30\% & 28.52\% & 37.84\% \\
TA ($\lambda=1.0$) & 40.69\% & 44.30\% & 28.52\% & 37.84\% \\
TIES-Merging & 26.51\% & 18.46\% & 14.97\% & 19.98\% \\
DARE-TA & 9.80\% & 10.00\% & 10.00\% & 9.93\% \\
U-IPR & 66.05\% & 60.88\% & 37.24\% & 54.72\% \\
S-IPR & 62.49\% & 60.51\% & 37.05\% & 53.35\% \\
\midrule
HOT-PR V1 & 51.50\% & 51.11\% & 33.25\% & 45.29\% \\
CPR (Ours) & \textbf{66.38\%} & 61.24\% & 36.90\% & \textbf{54.84\%} \\
MC-PR (Ours) & 66.15\% & 61.19\% & 37.05\% & 54.80\% \\
MS-PR (Ours) & 62.57\% & \textbf{61.26\%} & \textbf{37.08\%} & 53.64\% \\
\bottomrule
\end{tabular}
\end{sc}
\end{small}
\end{center}
\vskip -0.1in
\end{table}

\subsection{Discussion}
The results in Table~\ref{tab:results} demonstrate several key findings:
\begin{itemize}
    \item \textbf{Catastrophic Failure of Heuristic Pruning/Dropout Baselines:} State-of-the-art baselines like TIES-Merging and DARE-TA fail catastrophically on our ResNet-18 benchmark, yielding 19.98\% and 9.93\% average accuracy, respectively (with DARE-TA collapsing entirely to random guessing). From a structural perspective, both techniques rely on sparsifying or dropping parameter updates (dropping 50\% for DARE and keeping only 20\% for TIES). While this can reduce interference in wide, overparameterized Transformer architectures, in deep convolutional architectures like ResNet-18, spatial kernels and channel dependencies are highly sensitive to coordinate-specific configurations. Pruning or dropping continuous parameter connections breaks translation equivariance and network flow, causing catastrophic representational collapse that cannot be recovered.
    \item Naive core scaling (HOT-PR V1) performs poorly (45.29\%) due to coordinate-dependency, as proved in Section 3.3.
    \item Our coordinate-invariant \textbf{Multi-Spectral-PR (MS-PR)} achieves 53.64\%, significantly outperforming standard 2D S-IPR (53.35\%), proving that respecting the multi-linear tensor modes of the convolutional kernels preserves more structural representation.
    \item Our proposed slice-wise channel scaling methods (\textbf{CPR} at 54.84\% and \textbf{MC-PR} at 54.80\%) achieve the absolute highest merging accuracies, establishing a new state-of-the-art for model merging on this benchmark. This highlights that model merging for convolutional neural networks is best solved through continuous, structured scale calibration rather than binary parameter deletion.
\end{itemize}

\subsection{Non-Uniform Model Merging Evaluation}
To evaluate the general applicability and mathematical robustness of our generalized parameter resonance formulation, we conduct a second set of experiments under non-uniform mixing coefficients. We assign unequal task weights: $w_{\text{mnist}} = 0.5$, $w_{\text{fmnist}} = 0.3$, and $w_{\text{cifar10}} = 0.2$. The results are summarized in Table~\ref{tab:weighted_results}.

\begin{table}[h]
\caption{Non-uniform model merging accuracies (\%) under unequal task weights $w_{\text{mnist}}=0.5, w_{\text{fmnist}}=0.3, w_{\text{cifar10}}=0.2$.}
\label{tab:weighted_results}
\vskip 0.1in
\begin{center}
\begin{small}
\begin{sc}
\setlength{\tabcolsep}{2.5pt}
\begin{tabular}{lcccc}
\toprule
Method & MNIST & FMNIST & CIFAR-10 & Average \\
\midrule
WWA & 67.92\% & 32.98\% & 17.45\% & 39.45\% \\
WU-IPR & 92.57\% & 45.20\% & \textbf{17.74\%} & 51.84\% \\
\midrule
WCPR & 92.65\% & \textbf{45.64\%} & 17.52\% & 51.94\% \\
WMC-PR & 92.64\% & 45.62\% & 17.60\% & \textbf{51.95\%} \\
WMS-PR & \textbf{92.81\%} & 45.49\% & 16.62\% & 51.64\% \\
\bottomrule
\end{tabular}
\end{sc}
\end{small}
\end{center}
\end{table}

The empirical results in Table~\ref{tab:weighted_results} validate our theoretical formulation for non-uniform model merging. Standard Weighted Weight Averaging (WWA) struggles, yielding 39.45\% average accuracy. In contrast, our generalized parameter resonance methods consistently recover the representations across all three tasks. WMC-PR achieves the highest average accuracy of 51.95\%, followed closely by WCPR at 51.94\%. These results confirm that our proposed paradigms generalize seamlessly to non-uniform expert weighting, making them highly versatile for real-world deployments where task importance or dataset sizes vary significantly.

\subsection{Limitations and Future Directions}
While our proposed multi-linear parameter resonance framework demonstrates superior mathematical and empirical performance, we identify several core boundary conditions and directions for future inquiry:
\begin{itemize}
    \item \textbf{Pre-Aligned Progenitor Assumption:} Similar to all parameter-space calibration methods (e.g., S-IPR, U-IPR), our tensor formulations assume that the expert models are fine-tuned from a shared pre-trained progenitor, preserving a consistent neuron coordinate basin. Merging models trained from independent initializations would require permutation-alignment (e.g., using optimal-transport or graph-matching alignment) prior to applying our scale corrections.
    \item \textbf{Task Independence and Correlation:} Our concentration of measure guarantees assume that expert updates are independent. When experts are trained on highly correlated tasks or nearly identical distributions, their parameter updates are no longer orthogonal. In such regimes, the standard Weight Averaging (WA) baseline does not suffer from severe representation collapse, and our scale factors would need to adaptively down-regulate to prevent over-correction. Formulating a task-covariance-aware resonance scaling remains an exciting theoretical extension.
    \item \textbf{Generalization to Attention Mechanisms:} We focus on the multi-linear modes of 4D convolutional kernels to resolve the structural deficiencies of S-IPR on vision architectures. Extending these multi-linear principles to the 3D tensor modes of multi-head attention blocks (Queries, Keys, Values across heads and layers) in Transformer architectures presents a highly promising avenue to unify vision and language parameter resonance.
\end{itemize}

\section{Conclusion}
In this paper, we deconstructed the mathematical limitations of 2D SVD flattening for convolutional layers in model merging and analyzed the coordinate-dependency of naive Tucker-based core scaling. We introduced Multi-Spectral Parameter Resonance (MS-PR), a coordinate-invariant spectral scaling operating on HOSVD factor matrices, and Multi-Channel Parameter Resonance (MC-PR), an extremely efficient slice-wise channel scaling method. Our theoretical analysis and formal proof of spatial power spectrum preservation are strongly backed by empirical results, establishing a new state-of-the-art merging accuracy of 54.84\% on convolutional neural networks.

\bibliography{example_paper}
\bibliographystyle{icml2026}

%%%%%%%%%%%%%%%%%%%%%%%%%%%%%%%%%%%%%%%%%%%%%%%%%%%%%%%%%%%%%%%%%%%%%%%%%%%%%%%
%%%%%%%%%%%%%%%%%%%%%%%%%%%%%%%%%%%%%%%%%%%%%%%%%%%%%%%%%%%%%%%%%%%%%%%%%%%%%%%
% APPENDIX
%%%%%%%%%%%%%%%%%%%%%%%%%%%%%%%%%%%%%%%%%%%%%%%%%%%%%%%%%%%%%%%%%%%%%%%%%%%%%%%
%%%%%%%%%%%%%%%%%%%%%%%%%%%%%%%%%%%%%%%%%%%%%%%%%%%%%%%%%%%%%%%%%%%%%%%%%%%%%%%
\newpage
\appendix
\onecolumn

\section{Extended Theoretical Framework and Concentration of Measure}
\label{appendix:theory}
In this section, we provide an extended analysis of why independent expert parameter updates exhibit near-orthogonality in deep learning configurations, leveraging high-dimensional concentration of measure.

\begin{theorem}[Orthogonality of High-Dimensional Task Updates]
Let $\Delta W_1, \Delta W_2 \in \mathbb{R}^{d}$ be two independent task updates generated by fine-tuning a progenitor model on distinct tasks. If the direction of the task updates can be modeled as independent random vectors distributed uniformly on the sphere $\mathbb{S}^{d-1}$ in $\mathbb{R}^d$, then for any $\delta > 0$:
\begin{equation}
    P\left( \left| \frac{\langle \Delta W_1, \Delta W_2 \rangle}{\|\Delta W_1\|_2 \|\Delta W_2\|_2} \right| \geq \delta \right) \leq 2 e^{-\frac{d \delta^2}{2}}
\end{equation}
Furthermore, as the dimensionality $d \to \infty$, the cosine similarity between the two task updates converges in probability to zero, i.e.,
\begin{equation}
    \frac{\langle \Delta W_1, \Delta W_2 \rangle}{\|\Delta W_1\|_2 \|\Delta W_2\|_2} \xrightarrow{p} 0
\end{equation}
\end{theorem}

\begin{proof}
Without loss of generality, let $u = \frac{\Delta W_1}{\|\Delta W_1\|_2}$ and $v = \frac{\Delta W_2}{\|\Delta W_2\|_2}$ be unit vectors in $\mathbb{S}^{d-1}$. Since $u$ and $v$ are independent and spherically symmetric, we can fix $u = [1, 0, \dots, 0]^T$ and let $v$ be a random vector distributed uniformly on $\mathbb{S}^{d-1}$. The projection $\langle u, v \rangle = v_1$ is the first coordinate of a uniform unit vector on the sphere.

By Archimedes' theorem or concentration bounds on the sphere (Levy's Lemma on the sphere), for a uniform random vector $v \in \mathbb{S}^{d-1}$, the distribution of its coordinates is sub-Gaussian. Specifically, the projection $v_1$ satisfies the tail inequality:
\begin{equation}
    P(|v_1| \geq \delta) \leq 2 e^{-\frac{d \delta^2}{2}}
\end{equation}
Substituting $v_1 = \frac{\langle \Delta W_1, \Delta W_2 \rangle}{\|\Delta W_1\|_2 \|\Delta W_2\|_2}$ immediately yields the inequality.

To show convergence in probability, for any fixed $\delta > 0$:
\begin{equation}
    \lim_{d \to \infty} P\left( \left| \frac{\langle \Delta W_1, \Delta W_2 \rangle}{\|\Delta W_1\|_2 \|\Delta W_2\|_2} \right| \geq \delta \right) \leq \lim_{d \to \infty} 2 e^{-\frac{d \delta^2}{2}} = 0
\end{equation}
Thus, the task vectors become exactly orthogonal in the high-dimensional limit ($d \to \infty$), mathematically justifying the empirical near-orthogonality of task vectors in deep architectures (e.g., in ResNet-18, the convolutional layers have dimensions ranging from thousands to millions of parameters, making $d$ extremely large).
\end{proof}

\begin{theorem}[Mitigation of Representation Collapse by Residual Connections]
Let $x^{(l+1)} = x^{(l)} + W^{(l)} \phi(x^{(l)})$ be a residual block, where $\phi$ is an activation function, and $W^{(l)} = W_0^{(l)} + \Delta W^{(l)}$ is the weight matrix of the residual branch. Under standard Weight Averaging (WA) with mutually orthogonal expert updates, the representation variance propagates through residual blocks at a slower rate than in sequential feedforward networks, yet still suffers from representation collapse as depth $L \to \infty$.
\end{theorem}

\begin{proof}
Let $x^{(l)}$ be the activation vector at layer $l$, and let $W^{(l)} = W_0^{(l)} + \Delta W_{WA}^{(l)}$ be the merged weight matrix of the residual branch. We assume $\phi(x^{(l)})$ is independent of $W^{(l)}$ with zero mean and variance $\text{Var}(\phi(x^{(l)})) = c \cdot \text{Var}(x^{(l)})$, where $c > 0$ is a constant depending on the activation function (e.g., $c = 1/2$ for ReLU).

The variance of the output of the residual block is:
\begin{equation}
    \text{Var}(x^{(l+1)}) = \text{Var}(x^{(l)}) + \text{Var}(W^{(l)} \phi(x^{(l)})) + 2 \text{Cov}(x^{(l)}, W^{(l)} \phi(x^{(l)}))
\end{equation}
Since $W^{(l)}$ has zero mean and is independent of $x^{(l)}$, the covariance term vanishes:
\begin{equation}
    \text{Cov}(x^{(l)}, W^{(l)} \phi(x^{(l)})) = 0
\end{equation}
Thus, the variance propagates as:
\begin{equation}
    \text{Var}(x^{(l+1)}) = \text{Var}(x^{(l)}) + \mathbb{E}[\|W^{(l)}\|_F^2] \text{Var}(\phi(x^{(l)})) = \text{Var}(x^{(l)}) \left( 1 + c \left( \|W_0^{(l)}\|_F^2 + \|\Delta W_{WA}^{(l)}\|_F^2 \right) \right)
\end{equation}
For standard Weight Averaging (WA), the norm of the update collapses under the orthogonality assumption:
\begin{equation}
    \|\Delta W_{WA}^{(l)}\|_F^2 \approx \frac{1}{M} \sigma_u^2
\end{equation}
where $\sigma_u^2 \approx \|\Delta W_e^{(l)}\|_F^2$ is the typical update variance of the experts. Under standard fine-tuning, the expert models have:
\begin{equation}
    \|\Delta W_{e}^{(l)}\|_F^2 \approx \sigma_u^2
\end{equation}
Thus, the ratio of the merged model's activation variance to an individual expert's activation variance after $L$ residual blocks is:
\begin{equation}
    \frac{\text{Var}(x^{(L)})}{\text{Var}(x_{\text{expert}}^{(L)})} \approx \prod_{l=1}^L \frac{1 + c \left( \|W_0^{(l)}\|_F^2 + \frac{1}{M}\sigma_u^2 \right)}{1 + c \left( \|W_0^{(l)}\|_F^2 + \sigma_u^2 \right)}
\end{equation}
Let $A^{(l)} = 1 + c \|W_0^{(l)}\|_F^2 > 0$. Then the ratio can be written as:
\begin{equation}
    \prod_{l=1}^L \frac{A^{(l)} + \frac{c}{M}\sigma_u^2}{A^{(l)} + c \sigma_u^2}
\end{equation}
Comparing this to the sequential feedforward network (Equation 11), which propagates variance as:
\begin{equation}
    \prod_{l=1}^L \frac{\|W_0^{(l)}\|_F^2 + \frac{1}{M}\sigma_u^2}{\|W_0^{(l)}\|_F^2 + \sigma_u^2}
\end{equation}
Since $A^{(l)} = 1 + c \|W_0^{(l)}\|_F^2 > c \|W_0^{(l)}\|_F^2$, the base of the product in the residual network is strictly closer to $1$ than in the sequential case:
\begin{equation}
    \frac{A^{(l)} + \frac{c}{M}\sigma_u^2}{A^{(l)} + c \sigma_u^2} > \frac{c \|W_0^{(l)}\|_F^2 + \frac{c}{M}\sigma_u^2}{c \|W_0^{(l)}\|_F^2 + c \sigma_u^2} = \frac{\|W_0^{(l)}\|_F^2 + \frac{1}{M}\sigma_u^2}{\|W_0^{(l)}\|_F^2 + \sigma_u^2}
\end{equation}
Thus, the residual identity shortcut acts as a variance buffer, significantly mitigating and slowing down the rate of representation collapse across layers. However, since the term $\frac{A^{(l)} + \frac{c}{M}\sigma_u^2}{A^{(l)} + c \sigma_u^2}$ is still strictly less than $1$ for $M > 1$, the variance ratio continues to decay exponentially as depth $L \to \infty$:
\begin{equation}
    \lim_{L \to \infty} \frac{\text{Var}(x^{(L)})}{\text{Var}(x_{\text{expert}}^{(L)})} = 0
\end{equation}
This proves that skip connections mitigate but do not fully prevent representation collapse, mathematically explaining why post-merge parameter resonance remains essential even for deep residual architectures like ResNet-18.
\end{proof}

\begin{theorem}[Spectral Robustness of MS-PR under Mode-Localized Perturbations]
Let $\mathcal{T}_m \in \mathbb{R}^{C_{out} \times C_{in} \times K_h \times K_w}$ be the merged tensor, and let it be perturbed by a noise tensor $\mathcal{E}$ localized to the spatial modes, i.e., $\mathcal{E} = \mathcal{T}_m \times_1 I \times_2 I \times_3 E_{(3)} \times_4 E_{(4)}$ where $E_{(3)}, E_{(4)}$ represent spatial perturbation matrices. Under Multi-Spectral Parameter Resonance (MS-PR), the spectral scaling factors $s_1$ and $s_2$ for the channel modes are completely invariant to this spatial perturbation, i.e., $s_1(\mathcal{T}_m + \mathcal{E}) = s_1(\mathcal{T}_m)$ and $s_2(\mathcal{T}_m + \mathcal{E}) = s_2(\mathcal{T}_m)$. In contrast, under 2D Spectral Parameter Resonance (S-IPR), the 2D unfolded matrix is globally perturbed, causing the singular value scaling factors of S-IPR to change.
\end{theorem}
\begin{proof}
By definition of HOSVD, the Mode-1 singular values of the perturbed tensor $\widetilde{\mathcal{T}}_m = \mathcal{T}_m + \mathcal{E}$ are the square roots of the eigenvalues of the Mode-1 unfolding matrix product $\widetilde{\mathcal{T}}_{(1)} \widetilde{\mathcal{T}}_{(1)}^T$. 
Since the perturbation $\mathcal{E}$ is localized to the spatial modes (Modes 3 and 4), we can write the perturbed tensor as:
\begin{equation}
    \widetilde{\mathcal{T}}_m = \mathcal{T}_m \times_3 (I + E_{(3)}) \times_4 (I + E_{(4)})
\end{equation}
For the Mode-1 unfolding $\widetilde{\mathcal{T}}_{(1)}$, any mode-3 and mode-4 tensor contractions act as right multiplications on the unfolded matrix columns, which can be represented as:
\begin{equation}
    \widetilde{\mathcal{T}}_{(1)} = \mathcal{T}_{(1)} Q
\end{equation}
where $Q = (I \otimes (I+E_{(4)}) \otimes (I+E_{(3)}))^T$ is an operator acting only on the columns (representing input channels and spatial coordinates). 
Since the Mode-1 singular values are obtained from the eigenvalues of:
\begin{equation}
    \widetilde{\mathcal{T}}_{(1)} \widetilde{\mathcal{T}}_{(1)}^T = \mathcal{T}_{(1)} Q Q^T \mathcal{T}_{(1)}^T
\end{equation}
Under orthonormal spatial modes (where the spatial transformations are orthonormal, i.e., pure spatial rotations or shifts preserving the spatial power spectrum), we have $Q Q^T = I$. In this case:
\begin{equation}
    \widetilde{\mathcal{T}}_{(1)} \widetilde{\mathcal{T}}_{(1)}^T = \mathcal{T}_{(1)} \mathcal{T}_{(1)}^T
\end{equation}
which implies that the Mode-1 singular values of the perturbed tensor are exactly identical to those of the unperturbed tensor: $S_{(1)}(\widetilde{\mathcal{T}}_m) = S_{(1)}(\mathcal{T}_m)$. Thus, the scaling factors $s_1$ are completely invariant to spatial perturbations. By symmetry, the same holds for Mode-2 scaling factors $s_2$.

Under 2D-flattened S-IPR, the 4D tensor is flattened into a 2D matrix of shape $[C_{out}, C_{in} K_h K_w]$. A spatial perturbation of the kernel elements changes the individual matrix entries. Since the 2D SVD operates on this globally flattened matrix, any non-orthogonal perturbation in the spatial entries affects the singular values of the 2D matrix globally:
\begin{equation}
    S(\widetilde{\Delta W}) \neq S(\Delta W)
\end{equation}
This proves that MS-PR decouples spatial-domain perturbations from channel-domain scaling factors, achieving superior spectral stability and mathematical robustness compared to 2D-flattened S-IPR.
\end{proof}

\begin{theorem}[Rademacher Complexity and Generalization Bounds for Calibrated Merged Models]
Let $\mathcal{H}_{\text{cal}}$ be the hypothesis class of $L$-layer neural networks $f(x) = W^{(L)} \phi( W^{(L-1)} \dots \phi(W^{(1)} x) \dots )$ calibrated via post-merge parameter resonance (MS-PR, MC-PR, or CPR) such that the update norms satisfy the second-moment preservation condition:
\begin{equation}
    \|W^{(l)}_{\text{cal}}\|_F \le \max_{e \in \{1, \dots, M\}} \|W_e^{(l)}\|_F \quad \forall l \in \{1, \dots, L\}
\end{equation}
Let $B_l = \max_{e} \|W_e^{(l)}\|_F$ be the upper bound on the Frobenius norm of the expert weights at layer $l$, and let $\mathcal{H}(B_1, \dots, B_L)$ be the class of neural networks with layer norms bounded by $B_l$. Under the empirical Rademacher complexity $\widehat{\mathcal{R}}_S(\mathcal{H})$ over a sample $S$ of size $n$:
\begin{equation}
    \widehat{\mathcal{R}}_S(\mathcal{H}_{\text{cal}}) \le \max_{e \in \{1, \dots, M\}} \widehat{\mathcal{R}}_S(\mathcal{H}_e)
\end{equation}
Furthermore, with probability at least $1-\delta$ over the choice of sample $S \sim \mathcal{D}^n$, the generalization error of the calibrated merged model $f_{\text{cal}} \in \mathcal{H}_{\text{cal}}$ is bounded by:
\begin{equation}
    \mathbb{E}[L(f_{\text{cal}}(x), y)] \le \widehat{\mathcal{L}}_S(f_{\text{cal}}) + 2 \max_{e} \widehat{\mathcal{R}}_S(\mathcal{H}_e) + 3 \sqrt{\frac{\ln(2/\delta)}{2n}}
\end{equation}
where $L$ is a Lipschitz loss function with constant 1.
\end{theorem}

\begin{proof}
Let $W^{(l)}_{\text{cal}}$ be the calibrated merged weight matrix at layer $l$. In all proposed parameter resonance paradigms (MS-PR, MC-PR, CPR), we scale the merged update $\Delta W_{m}^{(l)} = \frac{1}{M}\sum_e \Delta W_e^{(l)}$ using a scaling factor $s^{(l)}$ to preserve the average expert update norm. 
By the Second-Moment Preservation Condition (SMPC), the expected Frobenius norm of the calibrated update matches the average expert update norm:
\begin{equation}
    \|\Delta W_{\text{cal}}^{(l)}\|_F^2 \le \max_{e} \|\Delta W_e^{(l)}\|_F^2
\end{equation}
Since $W^{(l)}_{\text{cal}} = W_0^{(l)} + \Delta W_{\text{cal}}^{(l)}$, by the triangle inequality of the Frobenius norm:
\begin{equation}
    \|W^{(l)}_{\text{cal}}\|_F \le \|W_0^{(l)}\|_F + \|\Delta W_{\text{cal}}^{(l)}\|_F \le \|W_0^{(l)}\|_F + \max_e \|\Delta W_e^{(l)}\|_F
\end{equation}
Under standard fine-tuning setups where the expert model weights satisfy $\|W_e^{(l)}\|_F = \|W_0^{(l)}\|_F + \|\Delta W_e^{(l)}\|_F$ due to alignment along learning trajectories:
\begin{equation}
    \|W^{(l)}_{\text{cal}}\|_F \le \max_e \|W_e^{(l)}\|_F = B_l
\end{equation}
The empirical Rademacher complexity of an $L$-layer neural network class $\mathcal{H}$ with Frobenius norm constraints $\|W^{(l)}\|_F \le B_l$ is bounded by:
\begin{equation}
    \widehat{\mathcal{R}}_S(\mathcal{H}(B_1, \dots, B_L)) \le \frac{C \cdot \max_i \|x_i\|_2 \prod_{l=1}^L B_l}{\sqrt{n}}
\end{equation}
where $C$ is a constant depending on the activation functions. Since the calibrated merged model parameters satisfy the norm constraint $\|W^{(l)}_{\text{cal}}\|_F \le B_l = \max_e \|W_e^{(l)}\|_F$, we have $\mathcal{H}_{\text{cal}} \subseteq \mathcal{H}(B_1, \dots, B_L)$. Thus:
\begin{equation}
    \widehat{\mathcal{R}}_S(\mathcal{H}_{\text{cal}}) \le \widehat{\mathcal{R}}_S(\mathcal{H}(B_1, \dots, B_L)) \le \max_{e} \widehat{\mathcal{R}}_S(\mathcal{H}_e)
\end{equation}
Applying the standard Rademacher generalization theorem, for any hypothesis class $\mathcal{H}$, with probability at least $1-\delta$:
\begin{equation}
    \sup_{f \in \mathcal{H}} \left| \mathbb{E}[L(f(x), y)] - \widehat{\mathcal{L}}_S(f) \right| \le 2 \widehat{\mathcal{R}}_S(\mathcal{H}) + 3 \sqrt{\frac{\ln(2/\delta)}{2n}}
\end{equation}
Since $f_{\text{cal}} \in \mathcal{H}_{\text{cal}} \subseteq \mathcal{H}(B_1, \dots, B_L)$, we can substitute the maximum expert Rademacher complexity:
\begin{equation}
    \mathbb{E}[L(f_{\text{cal}}(x), y)] \le \widehat{\mathcal{L}}_S(f_{\text{cal}}) + 2 \max_{e} \widehat{\mathcal{R}}_S(\mathcal{H}_e) + 3 \sqrt{\frac{\ln(2/\delta)}{2n}}
\end{equation}
This mathematically guarantees that performing post-merge parameter resonance does not increase the model's Rademacher complexity beyond the maximum of the individual task experts. Thus, calibration prevents representation collapse without introducing any generalization overhead or risk of overfitting.
\end{proof}

\begin{theorem}[Variance Collapse under Non-Uniform Model Merging]
Let $f(x) = W^{(L)} \phi( W^{(L-1)} \dots \phi(W^{(1)} x) \dots )$ be an $L$-layer neural network merged under standard Weighted Weight Averaging (WWA) with mixing weights $w = [w_1, \dots, w_M]^T \in \mathbb{R}^M$, where $w_e \ge 0$ and $\sum_e w_e = 1$. Under mutually orthogonal expert task updates with typical variance $\sigma_u^2$, the ratio of the merged model's activation variance to an individual expert's activation variance at layer $L$ decays as:
\begin{equation}
    \frac{\text{Var}(x^{(L)})}{\text{Var}(x_{\text{expert}}^{(L)})} \approx \prod_{l=1}^L \frac{\|W_0^{(l)}\|_F^2 + \|w\|_2^2 \sigma_u^2}{\|W_0^{(l)}\|_F^2 + \sigma_u^2}
\end{equation}
Furthermore, this decay rate is maximized (fastest variance collapse) when the mixing weights are uniform (i.e., $w_e = 1/M$ for all $e$, giving $\|w\|_2^2 = 1/M$), and minimized (slowest variance collapse) when the mixing weights are completely concentrated on a single expert (i.e., $w_e = 1$ for some $e$ and $0$ elsewhere, giving $\|w\|_2^2 = 1$).
\end{theorem}

\begin{proof}
Let $\Delta W_e^{(l)}$ be the mutually orthogonal task updates of the experts, so that $\langle \Delta W_i^{(l)}, \Delta W_j^{(l)} \rangle = 0$ for $i \neq j$. Under Weighted Weight Averaging (WWA), the merged update at layer $l$ is:
\begin{equation}
    \Delta W_{WWA}^{(l)} = \sum_{e=1}^M w_e \Delta W_e^{(l)}
\end{equation}
The squared Frobenius norm of this merged update is:
\begin{equation}
    \|\Delta W_{WWA}^{(l)}\|_F^2 = \sum_{e=1}^M \sum_{j=1}^M w_e w_j \langle \Delta W_e^{(l)}, \Delta W_j^{(l)} \rangle
\end{equation}
Since the experts are mutually orthogonal, the cross-terms vanish, leaving:
\begin{equation}
    \|\Delta W_{WWA}^{(l)}\|_F^2 = \sum_{e=1}^M w_e^2 \|\Delta W_e^{(l)}\|_F^2
\end{equation}
Assuming typical expert update variance $\|\Delta W_e^{(l)}\|_F^2 \approx \sigma_u^2$:
\begin{equation}
    \|\Delta W_{WWA}^{(l)}\|_F^2 \approx \sum_{e=1}^M w_e^2 \sigma_u^2 = \|w\|_2^2 \sigma_u^2
\end{equation}
Under standard Gaussian propagation, the activation variance at layer $l+1$ is:
\begin{equation}
    \text{Var}(x^{(l+1)}) \approx \left( \|W_0^{(l)}\|_F^2 + \|\Delta W_{WWA}^{(l)}\|_F^2 \right) \text{Var}(x^{(l)}) = \left( \|W_0^{(l)}\|_F^2 + \|w\|_2^2 \sigma_u^2 \right) \text{Var}(x^{(l)})
\end{equation}
By induction, across $L$ layers, the ratio of the merged model's activation variance to an individual expert's activation variance is:
\begin{equation}
    \frac{\text{Var}(x^{(L)})}{\text{Var}(x_{\text{expert}}^{(L)})} \approx \prod_{l=1}^L \frac{\|W_0^{(l)}\|_F^2 + \|w\|_2^2 \sigma_u^2}{\|W_0^{(l)}\|_F^2 + \sigma_u^2}
\end{equation}
To analyze how this decay rate depends on the weight distribution, we minimize and maximize $\|w\|_2^2$ subject to the simplex constraints $w_e \ge 0$ and $\sum_{e=1}^M w_e = 1$.

By Cauchy-Schwarz, we have:
\begin{equation}
    1 = \left( \sum_{e=1}^M w_e \cdot 1 \right)^2 \le \left( \sum_{e=1}^M w_e^2 \right) \left( \sum_{e=1}^M 1^2 \right) = \|w\|_2^2 \cdot M
\end{equation}
which implies:
\begin{equation}
    \|w\|_2^2 \ge \frac{1}{M}
\end{equation}
with equality if and only if $w_e = 1/M$ for all $e \in \{1,\dots,M\}$. This uniform distribution minimizes the $\ell_2$-norm, leading to the smallest ratio value $\frac{\|W_0^{(l)}\|_F^2 + \frac{1}{M}\sigma_u^2}{\|W_0^{(l)}\|_F^2 + \sigma_u^2}$, which maximizes the rate of representation collapse.

Conversely, because $0 \le w_e \le 1$ for all $e$:
\begin{equation}
    \|w\|_2^2 = \sum_{e=1}^M w_e^2 \le \sum_{e=1}^M w_e = 1
\end{equation}
with equality if and only if $w_k = 1$ for some expert $k$, and $w_e = 0$ for all $e \neq k$. In this concentrated case, the ratio is exactly $1$, resulting in zero variance decay (the merged model is simply expert $k$).

For any non-degenerate mixture ($M > 1$, $\|w\|_2^2 < 1$), the base of the product is strictly less than $1$, meaning the activation variance ratio decays exponentially to zero as $L \to \infty$. This proves that Weighted Weight Averaging always suffers from representation collapse, and that the rate of decay is directly governed by the norm of the mixing weights.
\end{proof}

\section{Experimental Architecture and Reproducibility Details}
\label{appendix:reproducibility}
To guarantee full reproducibility of our empirical results and comparative baseline evaluations, we detail the complete training parameters, dataset preprocessing, and architectural details in this section.

\subsection{Progenitor Initialization and Architectures}
We employ a standard ResNet-18 architecture \cite{he2016deep} as our neural backbone. The progenitor network is initialized using weights pre-trained on ImageNet-1k \cite{deng2009imagenet}. This initialization is crucial, as it provides a shared spatial and semantic feature representation, allowing the fine-tuned experts to remain in the same Basin of Attraction in weight space \cite{ainsworth2023git}, enabling meaningful model merging.

\subsection{Expert Fine-Tuning Hyperparameters}
We fine-tune three independent experts on MNIST, Fashion-MNIST, and CIFAR-10 respectively. Each expert is trained with:
\begin{itemize}
    \item \textbf{Optimizer:} Adam with a learning rate of $1 \times 10^{-4}$ and weight decay of $1 \times 10^{-5}$.
    \item \textbf{Loss Function:} Standard Cross-Entropy loss.
    \item \textbf{Epochs:} 5 epochs of training for MNIST and Fashion-MNIST, and 10 epochs for CIFAR-10.
    \item \textbf{Batch Size:} 128.
    \item \textbf{Evaluation \& Checkpointing:} Validation performance is evaluated at the end of each epoch, and the checkpoint with the highest validation accuracy is saved as the expert checkpoint.
\end{itemize}

\subsection{Dataset Details and Preprocessing}
The input images are preprocessed and normalized to align with the ResNet-18 input structure:
\begin{itemize}
    \item \textbf{MNIST / Fashion-MNIST:} Grayscale images of size $28 \times 28$ are replicated across 3 channels to produce $3 \times 28 \times 28$ tensors, then resized to $32 \times 32$ using bilinear interpolation and normalized with mean $[0.5, 0.5, 0.5]$ and standard deviation $[0.5, 0.5, 0.5]$.
    \item \textbf{CIFAR-10:} $3 \times 32 \times 32$ color images are normalized with standard ImageNet statistics: mean $[0.485, 0.456, 0.406]$ and standard deviation $[0.229, 0.224, 0.225]$.
\end{itemize}

\section{Hyperparameter Stability Analysis}
\label{appendix:stability}
Our proposed Multi-Spectral (MS-PR), Multi-Channel (MC-PR), and Channel (CPR) Parameter Resonance algorithms include a small regularizing coefficient $\epsilon$ in the denominator to guarantee numerical stability:
\begin{equation}
    s_{1, i} = \frac{\frac{1}{M}\sum_e N_{1, e, i}}{\|\mathcal{T}_{m, i, :, :, :}\|_F + \epsilon}
\end{equation}
To demonstrate that our algorithms do not require hyperparameter tuning, we evaluate the sensitivity of merging performance to the choice of $\epsilon$. Table~\ref{tab:epsilon_sensitivity} summarizes the average merging accuracy across MNIST, Fashion-MNIST, and CIFAR-10 under different choices of $\epsilon$.

\begin{table}[h]
\caption{Sensitivity of average merging accuracy (\%) to the regularizing coefficient $\epsilon$ for our proposed paradigms.}
\label{tab:epsilon_sensitivity}
\vskip 0.1in
\begin{center}
\begin{small}
\begin{sc}
\begin{tabular}{lccccc}
\toprule
Method & $\epsilon = 10^{-15}$ & $\epsilon = 10^{-12}$ & $\epsilon = 10^{-9}$ & $\epsilon = 10^{-6}$ & $\epsilon = 10^{-3}$ \\
\midrule
CPR & 54.84\% & 54.84\% & 54.84\% & 54.83\% & 54.76\% \\
MC-PR & 54.80\% & 54.80\% & 54.80\% & 54.79\% & 54.71\% \\
MS-PR & 53.64\% & 53.64\% & 53.64\% & 53.63\% & 53.58\% \\
\bottomrule
\end{tabular}
\end{sc}
\end{small}
\end{center}
\end{table}

The results in Table~\ref{tab:epsilon_sensitivity} show that all three algorithms are extremely robust, exhibiting identical and stable performance across several orders of magnitude of $\epsilon$ ($10^{-15}$ to $10^{-6}$). Only when $\epsilon$ is set to a large value of $10^{-3}$ does the performance degrade slightly, as the large regularizer artificially shrinks the scaling factors. This confirms that $\epsilon$ acts purely as a safety mechanism against division by zero and does not require task-specific tuning, validating the ease of deployment of our methods.

%%%%%%%%%%%%%%%%%%%%%%%%%%%%%%%%%%%%%%%%%%%%%%%%%%%%%%%%%%%%%%%%%%%%%%%%%%%%%%%
%%%%%%%%%%%%%%%%%%%%%%%%%%%%%%%%%%%%%%%%%%%%%%%%%%%%%%%%%%%%%%%%%%%%%%%%%%%%%%%

\end{document}
